\documentclass[10pt,conference]{IEEEtran}

\usepackage{amsmath}
\usepackage{booktabs}
\usepackage{graphicx}

\title{Edge-Enhanced Fashion-CNN: A Teaching-Oriented Baseline for Fashion-MNIST Classification}

\author{
\IEEEauthorblockN{Student Name}
\IEEEauthorblockA{Course Demo Project}
}

\begin{document}
\maketitle

\begin{abstract}
Fashion-MNIST is a widely used benchmark dataset for evaluating image classification algorithms on small-scale grayscale fashion images. In this demo paper, we reproduce a simple convolutional neural network baseline and introduce a teaching-oriented edge-enhanced variant. The proposed variant concatenates the original grayscale image with a Sobel edge map, allowing the model to access both appearance and contour information. We further include a random-noise control to examine whether the observed change is caused by meaningful edge information or merely by increasing the input dimensionality. This study is not intended to claim state-of-the-art performance; instead, it demonstrates a complete beginner-friendly research workflow, including baseline reproduction, minimal method modification, controlled comparison, and academic writing.
\end{abstract}

\section{Introduction}

Image classification is a fundamental task in machine learning and computer vision. Fashion-MNIST provides a compact yet more visually challenging alternative to the original MNIST dataset, as it contains grayscale images of fashion products from ten categories. Due to its small image size and standard train-test split, Fashion-MNIST is particularly suitable for beginner-level research training.

In this work, we use Fashion-MNIST as a case study to demonstrate how a beginner can move from baseline reproduction to a small method modification. We first implement a simple convolutional neural network as the baseline. Then, motivated by the observation that fashion categories often differ in object contours, we construct a Sobel edge map for each image and concatenate it with the original grayscale image as an additional input channel.

The goal of this work is not to propose a competitive model, but to provide a complete and reproducible research workflow. To avoid overclaiming, we additionally compare the edge-enhanced variant with a random-noise control, which helps examine whether the change comes from meaningful contour information or simply from increased input dimensionality.

\section{Method}

\subsection{Baseline CNN}

We adopt a lightweight convolutional neural network as the baseline model. The network contains two convolutional blocks followed by a fully connected classifier. Each convolutional block consists of a convolutional layer, batch normalization, ReLU activation, and max pooling.

\subsection{Edge-Enhanced Input}

Given an input image $X \in \mathbb{R}^{1 \times 28 \times 28}$, we compute its Sobel edge responses along the horizontal and vertical directions:
\begin{equation}
G_x = K_x * X, \quad G_y = K_y * X,
\end{equation}
where $K_x$ and $K_y$ denote the Sobel kernels and $*$ denotes convolution. The edge magnitude map is computed as:
\begin{equation}
E = \sqrt{G_x^2 + G_y^2}.
\end{equation}
The final input is constructed by concatenating the original image and the edge map:
\begin{equation}
\tilde{X} = \mathrm{Concat}(X, E).
\end{equation}
This changes the input from one channel to two channels while keeping the main CNN architecture unchanged.

\section{Experiments}

We compare three variants: a baseline CNN, an edge-enhanced CNN, and a random-noise control. The random-noise control is included to test whether any observed improvement comes from meaningful edge information or merely from adding a second input channel.

\begin{table}[t]
\centering
\caption{Comparison on Fashion-MNIST. Replace the placeholders with values from \texttt{results\_summary.csv}.}
\label{tab:results}
\begin{tabular}{llc}
\toprule
Method & Input Type & Test Accuracy \\
\midrule
Baseline CNN & Original image & xx.xx \\
Noise-Control CNN & Original + random noise & xx.xx \\
Edge-Enhanced CNN & Original + Sobel edge & xx.xx \\
\bottomrule
\end{tabular}
\end{table}

\section{Discussion}

The edge-enhanced variant provides a simple and interpretable modification to the input representation. If it outperforms the baseline and the noise-control variant, the result may suggest that contour information is useful for Fashion-MNIST classification. However, if the improvement is marginal or unstable, the result should be interpreted as a teaching-oriented demonstration rather than strong evidence of methodological superiority.

\section{Conclusion}

This demo project shows how to formulate a small research question, reproduce a baseline, implement a minimal method change, design a controlled comparison, and report results in a paper-like format.

\bibliographystyle{IEEEtran}
\bibliography{refs}

\end{document}
